\documentclass[literature.tex]{subfiles}

\begin{document}
\prose{Velký Gatsby}{Francis Scott Fitzgerald}
\teorieA{
epická próza, psychologický román (tragédie amerického snu)
}{
Modernismus, „Ztracená generace“ – kritika povrchnosti a materialismu „řvoucích dvacátých let“.
}{
Poetický, elegantní a ironický jazyk. Bohatá symbolika, přesné popisy, kontrast mezi lyrickými pasážemi a cynickým komentářem. Ich-forma vypravěče Nicka Carrawaye.
}{
Ironie, kontrast, opakování motivů (zelené světlo, večírky), symbolické detaily.
}{
Symbolismus (\texthl{zelené světlo = nedosažitelný americký sen}), metafora, alegorie (Gatsby jako ztělesnění amerického snu), ironie.
}{
Vyprávění v ich-formě z perspektivy Nicka Carrawaye s retrospektivními prvky. Elegantní, melancholický styl s lyrickými pasážemi a sarkastickým komentářem k vyšší společnosti. Krátké, intenzivní kapitoly.
}{
„Gatsby věřil v to zelené světlo; vzrušující budoucnost, která před námi rok od roku ustupuje. Dnes nám unikla, ale nevadí – zítra popoběhneme rychleji, rozpřáhneme paže dále... A jednoho krásného rána – 
A tak sebou zmítáme dál, lodě deroucí se proti proudu, bez přestání unášeni zpátky do minulosti.“
}
\teorieB{
\wtem{Jay Gatsby}{charismatický, záhadný milionář, který zbohatl nelegálně (pašování alkoholu). Idealista posedlý minulostí a láskou k Daisy. Symbol amerického snu.}
\wtem{Nick Carraway}{vypravěč, mladý muž ze Středozápadu, morální pozorovatel. Jediný, kdo zachová slušnost.}
\wtem{Daisy Buchananová}{krásná, povrchní, lehkomyslná žena z vyšší společnosti. Gatsbyho bývalá láska.}
\wtem{Tom Buchanan}{bohatý, arogantní, rasistický manžel Daisy, nevěrný a agresivní.}
\wtem{Jordan Bakerová}{profesionální golfistka, cynická, nemorální, Nickova přítelkyně.}
}{
Nick Carraway se přestěhuje na Long Island a stane se sousedem záhadného milionáře Jaye Gatsbyho. Gatsby pořádá velkolepé večírky, na které zve celou smetánku. Ukáže se, že Gatsby kdysi miloval Daisy (Nickovu sestřenici) a nyní se snaží získat ji zpět. Daisy je vdaná za bohatého, nevěrného Toma Buchanana. Po sérii setkání a hádek Daisy omylem srazí v Gatsbyho autě Tomovu milenku Myrtle Wilsonovou. Tom podstrčí vinu Gatsbymu. Myrtlein manžel Wilson Gatsbyho zastřelí a spáchá sebevraždu. Na Gatsbyho pohřeb přijde jen Nick a Gatsbyho otec. Daisy s Tomem odjedou a nechají vše za sebou.
}{
Chronologické vyprávění s retrospektivami (Gatsbyho minulost). Prostor: West Egg (nové bohatství) vs. East Egg (staré bohatství), New York. Čas: léto 1922 („řvoucí dvacátá léta“).
}{
Krutá kritika amerického snu, materialismu, povrchnosti a morálního úpadku vyšší společnosti. Peníze nezaručují štěstí ani lásku. Gatsbyho tragédie ukazuje, jak je minulost nedosažitelná (\texthl{zelené světlo} = iluze budoucnosti). „Lodě deroucí se proti proudu“ – lidé jsou neustále unášeni zpátky do minulosti.
}
\historie{
Období „řvoucích dvacátých let“ (Prohibition, jazz age, ekonomický růst před velkou hospodářskou krizí 1929).
}{}{}{}{
Klíčové dílo „Ztracené generace“. Spolu s Hemingwayem a Faulknerem patří k modernistické americké literatuře.
}{
(1896–1940). Americký spisovatel „Ztracené generace“. Proslavil se románem *This Side of Paradise*. Žil extravagantním životem s manželkou Zeldou. Alkoholismus a finanční problémy ho doprovázely celý život. Velký Gatsby byl za jeho života komerčním neúspěchem, uznání přišlo až posmrtně.
}{
Mezi nejznámější díla patří \textsc{Na tomto břehu ráje}, \textsc{Něžná je noc}, povídky jako \textsc{Bernard Q. a dívka s diamantovými očima}.
}\newpage
\kritika{}{}{
Kritika materialismu, falešných hodnot a iluzí amerického snu je stále velmi aktuální.
}
\nazor{
Velký Gatsby je mistrovské dílo – krátké, ale nesmírně hutné. Fitzgerald dokáže během 180 stran dokonale vykreslit prázdnotu bohaté společnosti. Nejsilnější je symbolika zeleného světla a tragický osud Gatsbyho, který celý život honil iluzi. Nick jako morální vypravěč je geniální volba. Kniha je smutná, ironická a krásně napsaná. Patří k vrcholům americké literatury 20. století – jedno z těch děl, které člověk přečte a už nikdy na něj nezapomene.
}
\explain{
\textbl{Americký sen --} ideál, že tvrdou prací a odhodláním může kdokoli dosáhnout úspěchu a štěstí;
\textbl{Ztracená generace --} spisovatelé, kteří prožili 1. světovou válku a ztratili iluze o světě;
\textbl{Symbolismus --} používání předmětů k vyjádření abstraktních myšlenek (zelené světlo).
}
\wpage
\end{document}